\subsubsection{Evolution of the Eigenspectrum}\label{sec:evolution-of-eigenspectrum}

Before discussing the eigenspectrum during annealing, let us recall a key result from \autopageref{sec:from-qubo-to-quantum-systems}, precisely from \autopageref{sec:the-eigenspectrum}. Given a Hamiltonian:
\begin{equation*}
    \hat{H}\ket{\psi_i} = E_i \ket{\psi_i}
\end{equation*}
The eigenvalues $E_i$ represent the possible energy levels of the system, while the eigenvectors $\ket{\psi_i}$ are the corresponding stationary states.

\highspace
\textcolor{Green3}{\faIcon{question-circle} \textbf{Why are stationary states important?}} We previously saw that if a system is in an eigenstate of the Hamiltonian:
\begin{equation*}
    \hat{H}\ket{\psi_i} = E_i \ket{\psi_i}
\end{equation*}
Then its evolution is given by:
\begin{equation*}
    \ket{\psi(t)} = e^{-i\hat{H}t/\hbar}\ket{\psi_i} = e^{-iE_{i}t/\hbar}\ket{\psi_i}
\end{equation*}
The state $\ket{\psi_i}$ only acquires a global phase factor $e^{-iE_{i}t/\hbar}$, which does not affect the physical properties of the state. Therefore, its measurement probabilities do not change over time. This is why we call $\ket{\psi_i}$ a \textit{stationary state}: it remains unchanged in terms of observable properties as time evolves.

\highspace
\begin{flushleft}
    \textcolor{Green3}{\faIcon{question-circle} \textbf{What changes during Quantum Annealing?}}
\end{flushleft}
During Quantum Annealing, the Hamiltonian of the system is time-dependent:
\begin{equation*}
    H(s) = A(s)H_{A} + B(s)H_{P}
\end{equation*}
As $s$ changes from $0$ to $1$, the \textbf{Hamiltonian evolves during the computation}. At the beginning ($s=0$), the system is in the ground state of $H_{A}$, which is easy to prepare. As $s$ increases, the Hamiltonian gradually transforms into $H_{P}$, whose ground state encodes the solution to our optimization problem. 

\highspace
Since the Hamiltonian changes with time, \hl{its eigenvalues also change}. Instead of having fixed energies $E_i$, we now have:
\begin{equation*}
    E_{0}(s) \leq E_{1}(s) \leq E_{2}(s) \leq \ldots
\end{equation*}
Which depend on the annealing parameter $s$. The \hl{collection of these energy levels is called the \textbf{eigenspectrum}}. The energy $E_i$ is the energy of the $i$-th eigenstate at the specific moment $s$ of the annealing process.

\highspace
It is crucial to understand that the \textbf{lowest energy level} is $E_{0}(s)$, which corresponds to the \textbf{instantaneous ground state}.
\begin{itemize}
    \item $E_{0}(s)$ is the ground-state energy;
    \item $E_{1}(s)$ is the energy of the first excited state;
    \item $E_{2}(s)$ is the energy of the second excited state;
    \item $\ldots$
\end{itemize}
So instead of a single energy landscape, we can imagine a family of landscapes that slowly deform as $s$ goes from $0$ to $1$.

\begin{figure}[!htp]
    \centering
    \begin{tikzpicture}[scale=1.1]

        % Axes
        \draw[->, thick] (0,0) -- (8,0) node[right] {$s$};
        \draw[->, thick] (0,0) -- (0,5) node[above] {Energy};

        % Labels on s axis
        \node[below] at (0,0) {$0$};
        \node[below] at (8,0) {$1$};

        % Energy curves
        \draw[very thick, Green3]
            plot[smooth] coordinates {
                (0,0.8) (1,0.9) (2,1.0) (3,1.15)
                (4,1.25) (5,1.15) (6,0.95) (7,0.75) (8,0.6)
            }
            node[right] {$E_0(s)$};

        \draw[very thick, orange]
            plot[smooth] coordinates {
                (0,2.0) (1,1.8) (2,1.55) (3,1.45)
                (4,1.65) (5,2.0) (6,2.4) (7,2.75) (8,3.0)
            }
            node[right] {$E_1(s)$};

        \draw[very thick, blue]
            plot[smooth] coordinates {
                (0,3.2) (1,3.0) (2,2.85) (3,2.7)
                (4,2.9) (5,3.25) (6,3.55) (7,3.8) (8,4.0)
            }
            node[right] {$E_2(s)$};

        % Gap indication
        \draw[<->, dashed, thick]
            (4,1.25) -- (4,1.65)
            node[midway, right] {$\Delta(s)$};

        % Arrow following ground state
        \draw[->, very thick, Green3]
            (1.2,0.55) .. controls (3,0.4) and (5,0.45) .. (6.8,0.55)
            node[midway, below] {follow ground state};

        % Text labels
        \node[align=center] at (1.5,4.5) {Initial\\Hamiltonian $H_A$};
        \node[align=center] at (6.5,4.5) {Problem\\Hamiltonian $H_P$};

        % Vertical guide lines
        \draw[dashed] (0,0) -- (0,4.3);
        \draw[dashed] (8,0) -- (8,4.3);

    \end{tikzpicture}
    \caption{Schematic representation of the evolution of the eigenspectrum during Quantum Annealing. The curves represent the energy levels $E_0(s)$, $E_1(s)$, and $E_2(s)$ as functions of the annealing parameter $s$. The gap $\Delta(s)$ between the ground state and the first excited state is called the \definition{Spectral Gap}, and it \textbf{measures how far the ground state is from the first excited state}. Gap is important because the smaller the gap becomes ($\Delta \approx 0$), the easier it is for the system to jump from the ground state to an excited state, which can lead to errors in the computation.}
\end{figure}